\documentclass{article}
\usepackage[utf8]{inputenc}
\usepackage[russian]{babel}
\usepackage{amsmath}
\usepackage{amsfonts}

\begin{document}

\section{Применение алгоритма MUSIC для оценки времени прихода (ToA)}

Алгоритм MUSIC (MUltiple SIgnal Classification) используется для высокоточной оценки времени прихода (ToA) в системах беспроводного позиционирования на основе стандарта IEEE 802.11ax. MUSIC позволяет эффективно разделять многолучевые компоненты сигнала, обеспечивая точную оценку временных задержек в сложных условиях распространения.

\subsection{Математическая модель}

Канальные оценки, полученные из HE-LTF (High Efficiency Long Training Field), представляют собой комплексную матрицу частотного отклика канала (CFR):
\begin{equation}
\mathbf{CFR} \in \mathbb{C}^{N \times P},
\end{equation}
где \( N \) --- количество активных поднесущих, \( P \) --- количество снимков (пространственно-временных потоков и приёмных антенн).

Для обеспечения равномерного частотного отклика пропущенные поднесущие интерполируются:
\begin{equation}
\mathbf{CFR}_{\text{full}}(k) = \text{interp1}(\text{activeFFTIndices}, \mathbf{CFR}, \text{allFFTIndices}),
\end{equation}
где интерполяция выполняется по амплитуде и фазе.

\subsection{Формирование корреляционной матрицы}

Корреляционная матрица формируется с применением пространственного сглаживания и усреднения по прямому и обратному направлениям для повышения устойчивости алгоритма MUSIC.

Для одного снимка (\( P = 1 \)) корреляционная матрица вычисляется как:
\begin{equation}
\mathbf{R} = \frac{1}{M} \sum_{k=1}^{M} \mathbf{y}_k \mathbf{y}_k^H,
\end{equation}
где \( \mathbf{y}_k = \mathbf{CFR}(k:k+L-1) \) --- подмассив длиной \( L = \lfloor \frac{2}{3} N \rfloor \), а \( M = N - L + 1 \) --- количество подмассивов.

Для нескольких снимков (\( P > 1 \)) сначала формируется выборочная корреляционная матрица:
\begin{equation}
\mathbf{R}_n = \frac{1}{P} \sum_{k=1}^{P} \mathbf{CFR}(:, k) \mathbf{CFR}(:, k)^H,
\end{equation}
а затем применяется пространственное сглаживание:
\begin{equation}
\mathbf{R} = \frac{1}{M} \sum_{k=1}^{M} \mathbf{R}_n(k:k+L-1, k:k+L-1).
\end{equation}

Для повышения устойчивости выполняется усреднение по прямому и обратному направлениям:
\begin{equation}
\mathbf{R} = \frac{\mathbf{R} + \mathbf{J} \mathbf{R}^* \mathbf{J}}{2},
\end{equation}
где \( \mathbf{J} \) --- матрица переворота, а \( \mathbf{R}^* \) --- комплексно-сопряжённая матрица.

\subsection{Псевдоспектр MUSIC}

Корреляционная матрица \( \mathbf{R} \) подвергается собственному разложению:
\begin{equation}
\mathbf{R} = \mathbf{V} \mathbf{\Lambda} \mathbf{V}^H,
\end{equation}
где \( \mathbf{E}_n \) --- матрица шумового подпространства, соответствующая \( L - M \) наименьшим собственным значениям (\( M \) --- число путей).

Псевдоспектр MUSIC вычисляется как:
\begin{equation}
P_{\text{MUSIC}}(\omega) = \frac{1}{\mathbf{a}(\omega)^H \mathbf{E}_n \mathbf{E}_n^H \mathbf{a}(\omega)},
\end{equation}
где \( \mathbf{a}(\omega) = [1, e^{j \omega}, \dots, e^{j (L-1) \omega}]^T \), а \( \omega = 2\pi \Delta f \tau \), \( \Delta f = \frac{\text{sampleRate}}{\text{fftLength}} \).

Временные задержки определяются из угловой частоты:
\begin{equation}
\tau = \frac{\omega}{2\pi \Delta f}.
\end{equation}

\subsection{Оценка ToA}

Пиковые значения \( P_{\text{MUSIC}}(\tau) \) соответствуют задержкам \( \tau_m \), где \( m = 1, 2, \dots, M \). Для исключения ложных пиков применяется порог:
\begin{equation}
\text{threshold} = \frac{\max(P_{\text{MUSIC}})}{10^{\text{dynamicRange}/10}},
\end{equation}
где \( \text{dynamicRange} \) зависит от ширины полосы канала.

ToA определяется как минимальная задержка:
\begin{equation}
\text{ToA} = \min(\tau_m).
\end{equation}

\subsection{Реализация в симуляции}

В реализованной симуляции алгоритм MUSIC применяется к канальным оценкам HE-LTF. Пространственное сглаживание и усреднение по прямому и обратному направлениям обеспечивают устойчивость в многолучевой среде. Псевдоспектр вычисляется с разрешением 0.1 м, что позволяет точно определять временные задержки. Полученное ToA используется для вычисления расстояний в алгоритмах трилатерации или комбинированного позиционирования с углами прихода (AoA).

\end{document}
